\chapter{Conclusion}

The thesis investigated whether symbolic chain-of-thought distillation (SCoTD) can improve a 3B-class small language model's math reasoning accuracy on GSM8K under strict greedy evaluation. It further compared rationale+answer supervision to answer-only (label-only) fine-tuning within a constrained single-GPU environment.

To address the specific research questions posed at the beginning of this thesis, the experimental results provide clear answers. Regarding \textbf{RQ1}, SCoTD improved the student from 70.51\% (base with CoT prompting) to 78.54\% (+8.03 points). Thus, supervised rationale transfer is effective for multi-step arithmetic in a 3B model under resource constraints. Regarding \textbf{RQ2}, label-only fine-tuning improved from 10.99\% to 15.01\% (+4.02 points) but remained far below SCoTD. This demonstrates that intermediate supervision is substantially more informative than final-answer-only supervision for this task and evaluation method.

The research conducted in this thesis led to the following key observations:
\begin{itemize}
    \item Distilling structured teacher reasoning traces yields significant accuracy gains without increased inference-time cost.
    \item Rationale diversity (30 generations/question) combined with correctness filtering and deduplication produces a high-signal training set.
    \item Parameter-efficient fine-tuning provides adequate adaptation capacity for reasoning improvements in small models under tight resource constraints.
    \item Strict formatting multiplies penalties for minor output deviations, especially for label-only training; evaluation design noticeably affects reported performance.
\end{itemize}

However, the study has limitations. The reliance on a single dataset and model, the absence of self-consistency decoding, and the lack of formatting error analysis limit the generalizability and depth of insights.

In conclusion, symbolic chain-of-thought distillation offers a cost-efficient method to elevate small model reasoning without reliance on large teacher inference at deployment time. Under constrained hardware, combining high-quality rationale curation with parameter-efficient fine-tuning yields meaningful accuracy gains and provides a reproducible path toward more capable small-scale reasoning systems.

\section{Future work}
Future work should address the identified gaps by extending the methodology to additional reasoning benchmarks, model sizes, and architectures. It would be beneficial to analyze the influence of formatting errors on evaluation, experiment with self-consistency decoding to measure potential performance ceilings, and test different distillation loss functions to further optimize knowledge transfer.
